\subsection{Two obstructions to diagonal-symmetry shortcuts}
\label{subsec:ff3-diagonal-shortcuts}

The next two observations explain why coprime complement conjugacy
does not settle the remaining full same-component boundary question.
They concern the entire Fermat geometric group, not just one
hyperelliptic quotient.

\begin{proposition}[A nontrivial diagonal twist cannot remove the residual obstruction]
\label{prop:ff3-diagonal-twist-obstruction}
Let $a\in\{5,7\}$, let $k=\mathbb F_Q$ have characteristic prime
to $3a$, and assume $\mu_a\subset k$.  On
$C_a:X^a+Y^a+Z^a=0$, let
\[
 \delta(X:Y:Z)=(\zeta^rX:\zeta^sY:Z),\qquad
 r,s\in\mathbb Z/a\mathbb Z,
\]
where $\zeta$ is primitive.  Let $N$ be the geometric pro-three
fundamental group and let $F$ denote arithmetic Frobenius.
If the action of $\delta^{-1}F$ on $N/\Phi(N)$ has three-power
order, then $r=s=0$.
In particular, for a nonidentity $\delta$, the closed subgroup of
$\operatorname{Out}(N)$ generated by its twisted Frobenius cannot
be pro-three.
\end{proposition}

\begin{proof}
Here $\delta$ and $F$ are initially outer automorphisms of $N$;
their actions on $N/\Phi(N)$ are canonical. They commute as outer
automorphisms because $\delta$ is defined over $k$. Thus the
displayed twisted action is defined without choosing a
$\delta$-fixed base point.

Put $J_a=\operatorname{Jac}(v^2+v=u^a)$, and write $\beta$ for
the automorphism induced by $(u,v)\mapsto(\zeta u,v)$.
The three degree-$a$ morphisms to this same curve are
\[
 \begin{split}
 f_Z&=(-XY/Z^2,X^a/Z^a),\\
 f_X&=(-YZ/X^2,Y^a/X^a),\\
 f_Y&=(-ZX/Y^2,Z^a/Y^a).
 \end{split}
\]
They satisfy $u^a=v(v+1)$ and extend to the proper smooth models.
Each is the quotient by its order-$a$ diagonal deck subgroup.
The induced pushforward on three-torsion is surjective: its
composition with pullback is multiplication by $a$, which is
invertible on three-torsion.

The three $u$-coordinate weights of $\delta$ are
\[
 w_Z=r+s,\qquad w_X=s-2r,\qquad w_Y=r-2s
 \pmod a.
\]
Consequently these quotient maps intertwine the action of
$\delta$ with $\beta^{w_i}$ and the action of $F$ with the same
Frobenius operator $B$ on $J_a[3]$.

The quotient of $v^2+v=u^a$ by $\beta$ is the projective $v$-line.
The $\beta$-invariants in its rational first homology therefore
vanish. Since $a$ is invertible in $\mathbb Z_3$, invariants are
an integral direct summand and commute with reduction modulo three.
Moreover $\dim_{\mathbb F_3}J_a[3]=a-1$, and $\Phi_a(T)$ is
irreducible over $\mathbb F_3$: the order of $3$ modulo $a$ is
$a-1$ for both $a=5$ and $a=7$. Thus $J_a[3]$ has rank one over
\[
 \mathbb F_3[\beta]\simeq\mathbb F_{3^{a-1}}.
\]
Frobenius commutes with $\beta$ because $\zeta\in k$, so $B$ is
a nonzero scalar in this field.

The three quotient actions of $\delta^{-1}F$ are therefore the
scalars $B\beta^{-w_i}$. Under the hypothesis each has
three-power order. Its order also divides $3^{a-1}-1$, which is
prime to three, so each is the identity. Hence
\[
 \beta^{w_Z}=\beta^{w_X}=\beta^{w_Y}.
\]
As $\beta$ has exact order $a$, this gives $3r=3s=0$ modulo $a$,
and therefore $r=s=0$. The final assertion follows because the
finite residual image of a pro-three group is a three-group.
\end{proof}

This excludes a nontrivial diagonal twist as the first step in the
pro-three-outer-action constant-extension construction. It does not
exclude other constructions of collisions, or decide the fibres
over a field where all boundary points are already rational.

\begin{proposition}[The geometric diagonal centralizer has a nonzero quotient]
\label{prop:ff3-diagonal-centralizer}
Let $a\in\{5,7\}$ and let $C_a$ be the proper Fermat curve over
an algebraically closed field of characteristic prime to $3a$.
Let $N$ be its geometric pro-three fundamental group and
$H=(\mu_a)^2$ its diagonal automorphism group. Form the relative
geometric action extension
\[
 1\longrightarrow N\longrightarrow E_H\longrightarrow H
 \longrightarrow1,
\]
retaining the full finite quotient $H$. Choose any complement
$S\subset E_H$, mapping isomorphically onto $H$, and let it act
on $N$ by conjugation. Such a complement exists by coprime
Schur--Zassenhaus. Put $L_2=\gamma_2N/\gamma_3N$ and
$g=(a-1)(a-2)/2$. Then
\[
 Z_N(S)\subset\gamma_2N,\qquad
 \frac{Z_N(S)}{Z_N(S)\cap\gamma_3N}
       \simeq L_2^H\simeq\mathbb Z_3^{g-1}.
\]
In particular this quotient has rank five for $a=5$ and rank
fourteen for $a=7$.
\end{proposition}

\begin{proof}
The choice of complement supplies an actual action on $N$.
Its action on each ordinary lower-central grade is the canonical
geometric $H$-action: changing lifts by inner automorphisms has
no effect on a grade. No arithmetic Galois invariance of $S$ is
assumed or asserted.

Write $V=N^{\mathrm{ab}}$. After an unramified coefficient
extension of $\mathbb Z_3$ containing $\mu_a$, the characters of
$H$ on $V$ are
\[
 (i,j)\in(\mathbb Z/a\mathbb Z)^2,
 \qquad i\ne0,\quad j\ne0,\quad i+j\ne0,
\]
each with multiplicity one. In characteristic zero, the
differentials $x^{i-1}y^{j-a}\,dx$, with $i,j\ge1$ and
$i+j\le a-1$, have characters $(i,j)$; adjoining their conjugates
gives precisely the displayed set. Smooth proper
prime-to-characteristic specialization gives the same geometric
representation in the stated characteristics. The degree-two
surface description used below is also given in
\citet[author pp.~2--4]{DavisPriesWickelgren2024}.
In particular $V^H=0$, so every element of $Z_N(S)$ lies in
$\gamma_2N$.

The degree-two surface presentation gives
\[
 L_2\simeq\bigwedge^2V/\mathbb Z_3\omega,
\]
where $\omega$ is the primitive symplectic relation. A trivial
diagonal character in the exterior square pairs $(i,j)$ with
$(-i,-j)$. There are exactly $g$ such pairs. The relation is
invariant and primitive, with a unit coefficient in a symplectic
basis. Since $|H|=a^2$ is invertible in $\mathbb Z_3$, the
averaging projector makes invariants exact. Therefore taking
invariants removes exactly one primitive rank-one summand,
and $L_2^H$ is finite free of rank $g-1$. The unramified
coefficient extension is faithfully flat, so this rank and
freeness assertion descends to $\mathbb Z_3$.

It remains to check that these graded fixed elements actually
lift to the geometric centralizer. For
$\overline x\in L_2^H$, choose $x\in\gamma_2N$.
The elements
\[
 c_h=x^{-1}h(x)\in\gamma_3N\qquad(h\in H)
\]
form a continuous nonabelian cocycle for the action furnished by
$S$. Its class is trivial because $H$ is finite of order prime
to three and $\gamma_3N$ is pro-three. Explicitly, finite
Schur--Zassenhaus conjugacy kills it in every finite
$H$-stable quotient; the corresponding nonempty closed sets of
conjugating elements have the finite-intersection property, so
compactness kills it in $\gamma_3N$ itself. Hence there is
$u\in\gamma_3N$ with $c_h=u\,h(u)^{-1}$ for every $h$.
Then $xu$ is fixed by $H$ and lifts $\overline x$.
Thus $Z_N(S)\to L_2^H$ is surjective, with kernel
$Z_N(S)\cap\gamma_3N$, proving the stated quotient.
\end{proof}

The absence of diagonal invariants in the original Jacobian does
not make this geometric centralizer trivial. Coprime conjugacy
of complements therefore does not provide the Frobenius
compatibility required for a rational invariant path. Neither
proposition proves a full same-component collision or full
separation, and neither supplies the missing affine
base-five or base-seven Selmer-section classification.
